% ID: SEL-D0033-OA-E04
% ESTADO: validada. S034, cotejada con examen y criterios.
\textbf{a)} $\hat p=74/121\approx0.61157$. Hay 74 casos
favorables y 47 restantes. Con aproximación normal bilateral
al $97\%$, $z_{0.985}\approx2.170090$:
\[
\hat p\pm z_{0.985}\sqrt{\frac{\hat p(1-\hat p)}{121}}
\approx\boxed{[0.51542,\ 0.70772]}.
\]
Los extremos del intervalo son aproximadamente $51.54\%$ y $70.77\%$.

\textbf{b)} Con la muestra fija, al bajar el nivel de confianza
disminuye el cuantil y, por tanto, el margen de error.

\textbf{c)} Las cuotas proporcionales son
$121/6$, $242/6$ y $363/6$, es decir, aproximadamente
$20.167$, $40.333$ y $60.5$. Al tomar las partes enteras suman 120.
Damos la plaza restante al tercer centro, que tiene el mayor
resto. La asignación es $\boxed{20,\ 40,\ 61}$, con total 121.

\textit{Asignación entera.} No existe un reparto entero exactamente
proporcional a $1:2:3$ con total 121, porque ese total no es múltiplo
de 6. Se ha utilizado el método de restos mayores para conservar el
tamaño total y aproximar las cuotas. El enunciado no fija un método
de redondeo y la asignación debe interpretarse con este criterio explícito.

